\documentclass[12pt,a4paper]{article}
\usepackage[russian]{babel}
\usepackage{fontspec}
\usepackage{geometry}
\usepackage{booktabs}
\usepackage{longtable}
\usepackage{array}
\usepackage{multirow}
\usepackage{enumitem}
\usepackage{titlesec}
\usepackage{xcolor}
\usepackage{siunitx}

\usepackage{pdfpages}

\usepackage{tablefootnote}

% Дополнительные пакеты для данного документа
\usepackage{caption}
\usepackage{subcaption}
\usepackage{listings}
% \usepackage{hyperref}

\geometry{a4paper, margin=25mm}
\setmainfont{Liberation Serif}
\setsansfont{Liberation Sans}

\definecolor{adblue}{RGB}{0,84,166}

\titleformat{\section}{\large\bfseries\color{adblue}}{}{0em}{}[\titlerule]
\titleformat{\subsection}{\normalsize\bfseries}{}{0em}{}
\titleformat{\subsubsection}{\normalsize\itshape}{}{0em}{}

\sisetup{output-decimal-marker={,}}

\usepackage{tikz}
\usetikzlibrary{positioning, calc, backgrounds, math, 3d, perspective, arrows.meta, graphs, graphdrawing, fit, shapes.geometric, trees}
    

\usepackage{pgfplots} % для построения графиков
\pgfplotsset{compat=1.18} % версия для совместимости
\usepackage{tkz-euclide}

% --- Подключение пакета алгоритмов (без флага algo2e, чтобы вернуть \begin{algorithm}) ---
\usepackage[ruled,vlined]{algorithm2e}
%\usepackage{caption}
\newcommand{\Gray}[1]{\textcolor{gray}{#1}} 
% Настройка подписей
\DeclareCaptionLabelFormat{gost}{#1 #2}
\captionsetup[table]{labelformat=gost,labelsep=endash,justification=justified,singlelinecheck=false,font=normal}
\captionsetup[figure]{labelformat=gost,labelsep=endash,justification=centering,font=normal}

% Локализация algorithm2e на русский (ОБЯЗАТЕЛЬНО после загрузки пакета)
\SetAlgorithmName{Алгоритм}{}{}
\SetAlgoCaptionSeparator{---}
\SetKwInput{Input}{Вход}
\SetKwInput{Output}{Выход}

\usepackage{url}
% Говорим пакету использовать моноширинный шрифт (как у texttt)
\Urlmuskip=0mu plus 1mu % гибкие пробелы для формул, если нужно
\DeclareUrlCommand\code{\urlstyle{tt}} 

% Библиотека схемотехники
\usepackage[siunitx, RPvoltages, american]{circuitikzgit}
\ctikzset{
    amplifiers/fill=cyan,
    sources/fill=green!20,
    diodes/fill=white,
    resistors/fill=violet,
    grounds/scale=1.2,
}

\usepackage{iftex}

\ifXeTeX % Если используется TeTeX или TeLaTeX
  \typeout{Режим XeTeX}
  \usepackage{xltxtra}
  \usepackage{fontspec}
  \usepackage{polyglossia}
  \setmainlanguage{russian}
  \setotherlanguage{english}
  \setkeys{russian}{babelshorthands=true} % По идее включает <<, >>, -- ---

  \setmainfont{Times New Roman}[Ligatures=TeX]
  \setromanfont{Times New Roman}[Ligatures=TeX]
  \setsansfont{Arial}[Ligatures=TeX]
  \setmonofont{Courier New}[Ligatures=TeX] 

  \newfontfamily{\cyrillicfont}{Times New Roman}
  \newfontfamily{\cyrillicfontrm}{Times New Roman}
  \newfontfamily{\cyrillicfonttt}{Courier New}
  \newfontfamily{\cyrillicfontsf}{Arial}
\fi

\ifLuaTeX % Если используется LuaLatex
  \typeout{Режим LuaTeX}
%  \usepackage{luatextra}
  \usepackage{fontspec}
  \defaultfontfeatures{Ligatures={TeX}}
  \usepackage{polyglossia}
  \setmainlanguage{russian}
  \setotherlanguage{english}
  \setmainfont{Times New Roman}[Script=Cyrillic, Ligatures=TeX]
  \setromanfont{Times New Roman}[Script=Cyrillic, Ligatures=TeX]
  \setsansfont{Arial}[Script=Cyrillic, Ligatures=TeX]
  \setmonofont{Courier New}[Ligatures=TeX]
  \newfontfamily{\cyrillicfont}{Times New Roman}[Script=Cyrillic, Ligatures=TeX]
  \newfontfamily{\cyrillicfontrm}{Times New Roman}[Script=Cyrillic, Ligatures=TeX]
  \newfontfamily{\cyrillicfonttt}{Courier New}
  \newfontfamily{\cyrillicfontsf}{Arial}[Script=Cyrillic, Ligatures=TeX]

  \usegdlibrary{trees}
  \usegdlibrary{force}
\fi

\ifPDFTeX % Если используется 8битный PDFTex
  \typeout{8-битный режим pdfTex}
  % For pdfTeX we must use old
  % 8-bit font technologies
  \usepackage[T2A]{fontenc}
  \usepackage[english,russian]{babel}
  \usepackage[utf8]{inputenc}
  \usepackage[english, russian]{babel}

  %\usepackage{newtxmath}
  \usepackage{amsmath}
  \usepackage{amsthm}
  \usepackage{amssymb}
  \usepackage{amsfonts}
  \usepackage{mathtext}
\fi

\usepackage{amsthm}  % Возможности AMSMath
\usepackage{amsmath}
\usepackage{amssymb}

\usepackage{esvect} % Для красивых стрелок
\usepackage{bm} % Для жирных символов

%Hyphenation rules
\usepackage{hyphenat}
\hyphenation{ма-те-ма-ти-ка вос-ста-нав-ли-вать ото-бра-жа-ют-ся сге-не-ри-ро-ван-ном сим-во-лы}
% \usepackage[english, russian]{babel}

\usepackage{graphicx}
\graphicspath{ {./images/} }

\usepackage{empheq}

\usepackage{float}

% \usepackage{comment}

\usepackage[hidelinks,
    unicode=true,          % поддерживает Unicode в закладках
    psdextra=true,      % автоматически конвертирует простую математику в текст
    % focusdefault=false
]{hyperref}


% Окружение "Теорема"
\newtheorem{theorem}{Теорема}[section]
\newtheorem{corollary}{Corollary}[theorem]
\newtheorem{lemma}[theorem]{Lemma}

% Окружение "Определение"
\theoremstyle{definition} % Прямой шрифт, нумерация
\newtheorem{definition}{Определение}[section]
\newtheorem{example}{Пример}[section]

% Окружение "Замечание"
\theoremstyle{remark} % Прямой шрифт, без нумерации
% Окуржение "Решение"
\newtheorem*{solution}{Решение}
\newtheorem*{remark}{Замечание}

\usepackage{polynom}
\usepackage{tabularx}
\usepackage{cancel}

\addto\captionsrussian{%
  \renewcommand{\figurename}{Рис.}%
  \renewcommand{\tablename}{Табл.}%
}

\renewcommand{\arraystretch}{1.5} % Высота строки таблицы

\renewcommand{\labelitemi}{\(\circ\)} % Бюллет -- пустой кружок в itemize

\title{Примечания к предварительной обработке}
\author{Наталия А. Рыбкина}
\date{2026}

\begin{document}

\maketitle

\section{Технология формирования обучающих пар: физическая модель деградации}
\label{sec:bayer_preparation}

В данном разделе описывается последовательность операций, которые превращают эталонное полноцветное изображение (High Quality, HQ) в реалистичный Low Quality (LQ) сигнал, соответствующий выходу малогабаритного CFA-сенсора.

\subsection{Порядок операций в конвейере деградации}

\begin{enumerate}
    \item \textbf{Пространственное сжатие (downscale)} – уменьшение разрешения в $k$ раз ($k=4$ или $k=8$). Моделирует удалённую сцену или мелкий размер пикселя.
    \item \textbf{Свёртка с функцией рассеяния точки (PSF)} – размытие, имитирующее частотно-контрастную характеристику (ЧКХ) объектива. Используется гауссово ядро с параметром $\sigma_{\text{PSF}}$.
    \item \textbf{Наложение маски Байера (паттерн RGGB)} – из полноцветного RGB-изображения создаётся одноканальная мозаика, где каждый пиксель сохраняет только один цветовой компонент (R, G$_1$, G$_2$ или B) согласно шахматной структуре.
    \item \textbf{Добавление шума} – генерируется аддитивный белый гауссов шум с заданным SNR, после чего он свёртывается с тем же ядром PSF (коррелированный шум) и прибавляется к мозаике. Так имитируется физический шум сенсора и усилителя.
    \item \textbf{Извлечение подканалов} – мозаичный сигнал раскладывается на четыре отдельных массива ($R$, $G_1$, $G_2$, $B$), которые сохраняются как 4-канальный 16-битный TIFF-файл (LQ-патч).
    \item \textbf{Сохранение эталона} – изображение, полученное после шагов 1–2 (сжатое и размытое, но без мозаики и шума), сохраняется как RGB-эталон (HQ) в формате PNG.
\end{enumerate}

\subsection{Почему маска Байера накладывается до добавления шума?}

В реальной камере свет сначала проходит через цветной фильтр (Байера), затем попадает на фотодиод, и только после этого к сигналу добавляются шумы квантования, тепловой шум и шум чтения. Таким образом, шум возникает \textbf{под тем же цветным фильтром}, который определил цвет пикселя. Если бы мы сначала зашумляли полноцветное изображение, а потом накладывали маску, то шум в разных цветовых каналах оказался бы независимым, что противоречит физике (зелёные пиксели шумят меньше, чем красные и синие). Поэтому порядок «Байер → шум» является единственно корректным.

\subsection{Почему аугментация ограничена поворотами на $90^\circ$ и отражениями?}

Структура маски Байера привязана к чётным и нечётным координатам пикселей. Поворот на $90^\circ$, $180^\circ$ или $270^\circ$ сохраняет эту чётность (просто меняет строки и столбцы местами), а поворот на произвольный угол (например, $30^\circ$) разрушает регулярную решётку. Следовательно, допустимы только преобразования из группы диэдра квадрата $D_4$ (8 вариантов). Использование любых других поворотов привело бы к появлению нереалистичных мозаичных паттернов, на которых нейросеть не сможет обобщаться на реальные данные.

\subsection{Что такое маска Байера и как она работает?}

Маска Байера — это не «разборка» на три отдельных изображения, а **одноканальная мозаика**, в которой каждый пиксель хранит значение только одного цвета (R, G или B). Например, для паттерна RGGB:

\[
\begin{array}{|c|c|c|c|}
\hline
R & G_1 & R & G_1 \\
\hline
G_2 & B & G_2 & B \\
\hline
R & G_1 & R & G_1 \\
\hline
G_2 & B & G_2 & B \\
\hline
\end{array}
\]

Зелёный цвет представлен двумя подканалами ($G_1$ и $G_2$), которые в сумме занимают половину площади матрицы. Красный и синий — по четверти. В результате получается **перемешанный** сигнал, по которому нельзя просто взять и получить три независимых полутоновых изображения. Задача демозаики (восстановления полного RGB) как раз и состоит в том, чтобы по этой мозаике корректно вычислить недостающие цвета в каждой точке.

\subsection{Краткие итоги}

\begin{itemize}
    \item Маска Байера **не разбирает** RGB на три слоя, а **кодирует** цвет в одной плоскости, теряя два других канала в каждом пикселе.
    \item Порядок обработки: сжатие → размытие → Байер → шум → извлечение четырёх каналов.
    \item Аугментация допускает только повороты на $90^\circ$ и отражения (группа $D_4$).
    \item Такой конвейер максимально приближает синтетические данные к реальному RAW-сигналу малогабаритного сенсора.
\end{itemize}

\end{document}